% GENERATED FILE. Edit canonical lessons, then run npm run generate:books.
% canonical-source-hash: fnv1a64:e988f5b235001ba1
% canonical-lessons: IT-C525-libraio, IT-C525-copertina, IT-C525-sfogliare, IT-C525-segnalibro, IT-C525-tascabile

\chapter{In libreria}
\label{ch:it-in-libreria}
\hlchaptermodality{525}

\begin{chapteropening}
\textbf{By the end of this chapter:} \emph{I can talk about choosing and handling books in a bookshop.}

Complete the last lesson of chapter 525: I can talk about choosing and handling books in a bookshop.
\end{chapteropening}

\section[il libraio]{il libraio — bookseller}

\label{lesson:IT-C525-libraio}



\begin{quote}
Before the new one: say the Italian for twist, then the Italian for ending.
\end{quote}

\subsection*{You'll want to know: il libraio}
\textbf{il libraio} — \textquotedblleft{}bookseller\textquotedblright{}.

\begin{quote}
Il libraio mi ha consigliato un romanzo giallo.
\end{quote}

\begin{quote}
\emph{The bookseller recommended a crime novel to me.}
\end{quote}

\emph{La libreria} is the bookshop; \emph{la biblioteca} is the library.

\begin{cousinweb}
\emph{libraio} is built on \emph{libro} \textquotedblleft{}book\textquotedblright{}, inherited from Latin \emph{liber} \textquotedblleft{}the inner bark of a tree, a book\textquotedblright{} (Lewis \& Short).
\end{cousinweb}

\begin{grammarlens}[title={What you've built}]
One more word for telling, explaining and describing at B1.
\end{grammarlens}

\subsection*{Your turn}
\begin{itemize}
\raggedright
  \item \emph{Say it:} \emph{il libraio}
  \item \emph{Say it:} \emph{il libraio}, once more
  \item \emph{Say it:} a sentence of your own with \emph{il libraio}
  \item \emph{Recall:} say the Italian for twist, then the Italian for ending, then say \emph{il libraio} again
\end{itemize}

\begin{tcolorbox}[breakable,colback=teal!4,colframe=teal!35!black,title={Before you move on}]
What does il libraio mean? (\textquotedblleft{}Bookseller\textquotedblright{}.) Say it once more.
\end{tcolorbox}
\section[la copertina]{la copertina — cover (of a book)}

\label{lesson:IT-C525-copertina}



\begin{quote}
Before the new one: say the Italian for ending, then the Italian for bookseller.
\end{quote}

\subsection*{You'll want to know: la copertina}
\textbf{la copertina} — \textquotedblleft{}cover (of a book)\textquotedblright{}.

\begin{quote}
Mi piace la copertina di questo libro.
\end{quote}

\begin{quote}
\emph{I like the cover of this book.}
\end{quote}

Built on \emph{coperta} \textquotedblleft{}blanket, cover\textquotedblright{}.

\begin{cousinweb}
\emph{copertina} is a little \emph{coperta} \textquotedblleft{}cover\textquotedblright{}, from \emph{coprire} \textquotedblleft{}to cover\textquotedblright{}, inherited from Latin \emph{cooperīre} \textquotedblleft{}to cover over\textquotedblright{} (Lewis \& Short).
\end{cousinweb}

\begin{grammarlens}[title={What you've built}]
One more word for telling, explaining and describing at B1.
\end{grammarlens}

\subsection*{Your turn}
\begin{itemize}
\raggedright
  \item \emph{Say it:} \emph{la copertina}
  \item \emph{Say it:} \emph{la copertina}, once more
  \item \emph{Say it:} a sentence of your own with \emph{la copertina}
  \item \emph{Recall:} say the Italian for ending, then the Italian for bookseller, then say \emph{la copertina} again
\end{itemize}

\begin{tcolorbox}[breakable,colback=teal!4,colframe=teal!35!black,title={Before you move on}]
What does la copertina mean? (\textquotedblleft{}Cover (of a book)\textquotedblright{}.) Say it once more.
\end{tcolorbox}
\section[sfogliare]{sfogliare — to leaf through}

\label{lesson:IT-C525-sfogliare}



\begin{quote}
Before the new one: say the Italian for bookseller, then the Italian for cover.
\end{quote}

\subsection*{You'll want to know: sfogliare}
\textbf{sfogliare} — \textquotedblleft{}to leaf through\textquotedblright{}.

\begin{quote}
Sfoglio una rivista mentre aspetto.
\end{quote}

\begin{quote}
\emph{I leaf through a magazine while I wait.}
\end{quote}

Built on \emph{foglio} \textquotedblleft{}sheet\textquotedblright{}: turning the sheets.

\begin{cousinweb}
\emph{sfogliare} is built from \emph{s-} and \emph{foglio} \textquotedblleft{}sheet of paper\textquotedblright{}, inherited from Latin \emph{folium} \textquotedblleft{}a leaf\textquotedblright{} (Lewis \& Short).
\end{cousinweb}

\begin{grammarlens}[title={What you've built}]
One more word for telling, explaining and describing at B1.
\end{grammarlens}

\subsection*{Your turn}
\begin{itemize}
\raggedright
  \item \emph{Say it:} \emph{sfogliare}
  \item \emph{Say it:} \emph{sfogliare}, once more
  \item \emph{Say it:} a sentence of your own with \emph{sfogliare}
  \item \emph{Recall:} say the Italian for bookseller, then the Italian for cover, then say \emph{sfogliare} again
\end{itemize}

\begin{tcolorbox}[breakable,colback=teal!4,colframe=teal!35!black,title={Before you move on}]
What does sfogliare mean? (\textquotedblleft{}To leaf through\textquotedblright{}.) Say it once more.
\end{tcolorbox}
\section[il segnalibro]{il segnalibro — bookmark}

\label{lesson:IT-C525-segnalibro}



\begin{quote}
Before the new one: say the Italian for cover, then the Italian for to leaf through.
\end{quote}

\subsection*{You'll want to know: il segnalibro}
\textbf{il segnalibro} — \textquotedblleft{}bookmark\textquotedblright{}.

\begin{quote}
Ho perso il segnalibro e non so a che pagina ero.
\end{quote}

\begin{quote}
\emph{I lost my bookmark and do not know what page I was on.}
\end{quote}

From \emph{segnare} and \emph{libro}.

\begin{cousinweb}
\emph{segnalibro} is built from \emph{segna} \textquotedblleft{}marks\textquotedblright{}, from Latin \emph{signāre} (EtymDB), and \emph{libro} \textquotedblleft{}book\textquotedblright{}, from Latin \emph{liber} (Lewis \& Short).
\end{cousinweb}

\begin{grammarlens}[title={What you've built}]
One more word for telling, explaining and describing at B1.
\end{grammarlens}

\subsection*{Your turn}
\begin{itemize}
\raggedright
  \item \emph{Say it:} \emph{il segnalibro}
  \item \emph{Say it:} \emph{il segnalibro}, once more
  \item \emph{Say it:} a sentence of your own with \emph{il segnalibro}
  \item \emph{Recall:} say the Italian for cover, then the Italian for to leaf through, then say \emph{il segnalibro} again
\end{itemize}

\begin{tcolorbox}[breakable,colback=teal!4,colframe=teal!35!black,title={Before you move on}]
What does il segnalibro mean? (\textquotedblleft{}Bookmark\textquotedblright{}.) Say it once more.
\end{tcolorbox}
\section[il tascabile]{il tascabile — paperback; pocket-sized}

\label{lesson:IT-C525-tascabile}



\begin{quote}
Before the new one: say the Italian for to leaf through, then the Italian for bookmark.
\end{quote}

\subsection*{You'll want to know: il tascabile}
\textbf{il tascabile} — \textquotedblleft{}paperback; pocket-sized\textquotedblright{}.

\begin{quote}
Compro i tascabili perché costano meno.
\end{quote}

\begin{quote}
\emph{I buy paperbacks because they cost less.}
\end{quote}

As an adjective: \emph{un'edizione tascabile}.

\begin{cousinweb}
\emph{tascabile} is built on \emph{tasca} \textquotedblleft{}pocket\textquotedblright{}, which EtymDB traces to Frankish \emph{taska} \textquotedblleft{}pouch\textquotedblright{}; others start from Late Latin \emph{tasca}, so the origin is not fully settled.
\end{cousinweb}

\begin{grammarlens}[title={What you've built}]
One more word for telling, explaining and describing at B1.
\end{grammarlens}

\subsection*{Your turn}
\begin{itemize}
\raggedright
  \item \emph{Say it:} \emph{il tascabile}
  \item \emph{Say it:} \emph{il tascabile}, once more
  \item \emph{Say it:} a sentence of your own with \emph{il tascabile}
  \item \emph{Recall:} say the Italian for to leaf through, then the Italian for bookmark, then say \emph{il tascabile} again
\end{itemize}

\begin{tcolorbox}[breakable,colback=teal!4,colframe=teal!35!black,title={Before you move on}]
What does il tascabile mean? (\textquotedblleft{}Paperback; pocket-sized\textquotedblright{}.) Say it once more.
\end{tcolorbox}
